% ==============================================================================
% ПАРТИЯ B03-006: Глава 1: Статистика функционирования и базовые проблемы РНМ
% Диапазон: DISS-005-B0236 .. DISS-005-B0259 (24 блока)
% ==============================================================================

% DISS-005-B0236
\section{Статистика функционирования\commentnote{DISS-005-C0009}{Евдаков Алексей}{Пока почти полностью отсутствует какой-либо анализ.} БАВР}\label{sec:ch1_statistics}

% DISS-005-B0237
Появится описание статистики. Есть идеи по следующим задачам\par\vspace{0.2cm}

% DISS-005-B0238
- Статистика на симулированных данных (то чем сейчас занимается и будет заниматься Данил)\par\vspace{0.2cm}

% DISS-005-B0239
- Статистика по работе БАВР реальных БАВР. Пока сложно сказать, как это сделать, можно прост описать задачу и на этом поднять актуальность сбора данных.\par\vspace{0.2cm}

% DISS-005-B0240
- Статистика по выключателями (данные есть, надо ток написать код, это не сложно). Чтобы было понятно, какие реальные времена выключателей.\par\vspace{0.2cm}

% DISS-005-B0241
\par\vspace{0.2cm}
% DISS-005-B0242
\noindent\textbf{(Проблема\commentnote{DISS-005-C0010}{Евдаков Алексей}{См. файл D:\textbackslash{}Буфер (Алексей)\textbackslash{}Отчёты\textbackslash{}Саратов НПЗ (Роснефть)\textbackslash{}29.05.2025. Там отправили запрос – как справиться с срабатыванием при синхронной просадке?}) Синхронная просадка напряжения:}\par\vspace{0.2cm}

% DISS-005-B0243
Часто говорят о такой проблеме, но не существует строго определения применимого для функционирования БАВР. Надо изучить этот вопрос и принять какое-то решение. Потому что кто-то говорит, что было синхронное, а кто-то не соглашается, ведь от чего считать эту синхронность?\par\vspace{0.2cm}

% DISS-005-B0244
\par\vspace{0.2cm}
% DISS-005-B0245
\noindent\textbf{ВЫВОДЫ:}\par\vspace{0.2cm}
% DISS-005-B0246
Из данной части надо привести анализ сравнения различных производителей. Плюс сделать акцент на том, что часто спорным органом является РНМ, и поэтому он требует более подробного рассмотрения\par\vspace{0.3cm}

% DISS-005-B0247 .. DISS-005-B0249
\par\vspace{0.3cm}
% DISS-005-B0250
\section{Базовые проблемы функционирования РНМ}\label{sec:ch1_pdr_problems}

% DISS-005-B0251
Существует 2 подхода к оценке направления мощности – это оценка по каждой из фаз, и оценка суммарной мощности. Пофазный РНМ принято рассчитывать на основе тока указанной фазы и линейного напряжения двух других фаз. Суммарный РНМ в большинстве случаев рассчитывается на основе тока и напряжения прямой последовательности. Хотя могут использоваться и проекции активной / реактивной мощности на различные оси, но такие случае встречаются не часто.\par\vspace{0.2cm}

% DISS-005-B0252
Для БАВР пофазную оценку рекомендуется использовать в сетях с заземлённой нейтралью, так как там, при аварии на одной из фаз влияние на две другие может не оказываться, и в итоге использование органов на суммарной мощности может привести к отсутствию срабатывания, что как было показано в предыдущей главе (Раздел ХХ\commentnote{DISS-005-C0011}{Евдаков Алексей}{Вставить ссылку}).\par\vspace{0.2cm}

% DISS-005-B0253
\noindent\commentnote{DISS-005-C0012}{Feel}{Если публиковать, то с расшифровкой каждой буквы + советую нарисовать подключение БАВР и обозначить ВСЕ величины, как в формулах/на рисунках. Везде одинаково. Думаю, лучше будет смотреться в диссере, в тексте. Там в любом случае надо будет описать. А вы формулы патентуете?}\commentnote{DISS-005-C0013}{Евдаков Алексей (ответ на C0012)}{Нет, не патентуем. Поэтому пока и вопрос о разглашении формулы}\commentnote{DISS-005-C0014}{Евдаков Алексей (ответ на C0012)}{Да и в целом, здесь не сказать, что прям очень нужно указывать. Может кратко будет описано ранее в работах Данила}\par\vspace{0.2cm}

% DISS-005-B0254
\noindent\textit{В РЭ нет формул, требуется согласовать их указание здесь (есть в ПРО версии ПО просмотра осциллограмм). У других производителей формула отсутствует.}\par\vspace{0.2cm}

% DISS-005-B0255
Реле направления мощности основанный на прямой последовательности (РНМПП) применяется в сетях с изолированной нейтралью, так как в них аварийными считаются только 2-ух и 3-ёх фазные замыкания. При этом, РНМПП обладает значительно большей устойчивостью и не происходит излишних возвратов органа при внешних замыканиях.\par\vspace{0.2cm}

% DISS-005-B0256
\par\vspace{0.2cm}

% DISS-005-B0257
\noindent\textit{В РЭ нет формул, требуется согласовать их указание здесь (есть в ПРО версии ПО просмотра осциллограмм). У других производителей формула отсутствует.}\par\vspace{0.2cm}

% DISS-005-B0258
\par\vspace{0.2cm}

% DISS-005-B0259
Так как большинство БАВР установлены в сетях 6-35кВ, а их внедрение в сети 0,4кВ (сети с заземлённой нейтралью), только начинает набирать массовость, то дальше все проблемы описаны на основе РНМПП\par\vspace{0.2cm}
